\documentclass{training}
\schoollogo{../../resources/intros/UI-02-rules.jpg}
\newcommand{\uxfond}{../../resources/fond/UI-02-rules.png}
% Filigrane PAR MODULE : \uxfond (PNG A4 dilue 30 %). Cadre noir vectoriel commun.
% Proportions original 744 px : marge 20/744, trait 1/744. Exterieur BLANC.
% Cadre via eso-pic +\put (coordonnees absolues) — pas de remember picture (decalage).
\usepackage{eso-pic}

\providecommand{\uxfond}{../../resources/fond/norme_fond_paliers_30.png}

\newlength{\uxmargin}
\newlength{\uxthick}
\newlength{\uxinnerw}
\newlength{\uxinnerh}
\setlength{\uxmargin}{0.02688\paperwidth}% 20/744
\setlength{\uxthick}{0.001344\paperwidth}% 1/744
\setlength{\uxinnerw}{\dimexpr\paperwidth-2\uxmargin-2\uxthick\relax}
\setlength{\uxinnerh}{\dimexpr\paperheight-2\uxmargin-2\uxthick\relax}

% Contenu strictement a l'interieur du trait (+ jour confortable).
\geometry{%
  includeheadfoot,
  left=22mm,
  right=22mm,
  top=20mm,
  bottom=18mm,
  headheight=30pt,
  headsep=5mm,
  footskip=12mm%
}

\AddToShipoutPictureBG{%
  \AtPageLowerLeft{%
    % Fond blanc page entiere
    \put(0,0){\color{white}\rule{\paperwidth}{\paperheight}}%
    % Filigrane module a l'interieur du cadre
    \put(\LenToUnit{\dimexpr\uxmargin+\uxthick\relax},\LenToUnit{\dimexpr\uxmargin+\uxthick\relax}){%
      \includegraphics[width=\uxinnerw,height=\uxinnerh]{\uxfond}}%
    % Cadre noir : bas, haut, gauche, droite
    \put(\LenToUnit{\uxmargin},\LenToUnit{\uxmargin}){%
      \color{black}\rule{\dimexpr\paperwidth-2\uxmargin\relax}{\uxthick}}%
    \put(\LenToUnit{\uxmargin},\LenToUnit{\dimexpr\paperheight-\uxmargin-\uxthick\relax}){%
      \color{black}\rule{\dimexpr\paperwidth-2\uxmargin\relax}{\uxthick}}%
    \put(\LenToUnit{\uxmargin},\LenToUnit{\uxmargin}){%
      \color{black}\rule{\uxthick}{\dimexpr\paperheight-2\uxmargin\relax}}%
    \put(\LenToUnit{\dimexpr\paperwidth-\uxmargin-\uxthick\relax},\LenToUnit{\uxmargin}){%
      \color{black}\rule{\uxthick}{\dimexpr\paperheight-2\uxmargin\relax}}%
  }%
}

\makeatletter
\fancyhf{}
\fancyheadoffset{0pt}
\renewcommand{\headrulewidth}{0pt}
\lhead{\parbox[t]{0.56\textwidth}{\RaggedRight\footnotesize\sloppy\@title}}
\rhead{\parbox[t]{0.36\textwidth}{\RaggedLeft\footnotesize\sloppy\@subtitle}}
\cfoot{\thepage}
% Coupures : underscore et slash dans les noms de fichiers / consignes
\DeclareRobustCommand{\_}{\textunderscore\allowbreak}
\makeatother

% Evite les depassements de longues lignes (Github, Name_FirstName_...)
\sloppy
\emergencystretch=3em
\hbadness=10000

% Sommaire type 42 : chapitres seulement
\setcounter{tocdepth}{0}
\renewcommand{\cftchappresnum}{}
\renewcommand{\cftchapaftersnum}{}
\setlength{\cftchapnumwidth}{2.0em}
\setlength{\cftbeforechapskip}{0.35em}
\cftsetindents{chapter}{0em}{2.0em}


\begin{document}

\title{Rules}
\subtitle{Piscine UX/UI --- UI}
\member{Piscine UX/UI}{https://github.com/01-edu/UX-UI}
\summary{Grids, states, accessibility, and the small controls people reuse.}
\version{1.0}
\prerequisite{The previous module of this piscine, or none if this is your first.}
\maketitle

\newpage
\tableofcontents
\newpage


\chapter{Instructions}

The foreword is not an exercise. It will not be graded. The exercises start at the next chapter, from the easiest contract to the hardest. A later exercise is not taken into account if an earlier one does not hold.

Only what is in the turn-in directory of an exercise is looked at. Do not add files that the exercise does not name. The name of a deliverable follows the pattern in the exercise: \texttt{Name\_FirstName\_DeliverableName\_Date\_VersionNumber}.

You will not be given the gesture. The examples in an exercise sometimes carry a constraint that the sentence does not repeat. Read them. Your other references are the manual of the tool you picked, the person sitting next to you, and the links printed in the exercise.

The order is the order of the quest. Skipping a step because a later screen looks more interesting produces a screen nobody can explain.

\subsection*{About the pictures}

The mark on the cover is the picture of this module. The pale steps behind the text are the page ground. Neither one is a hint.

\subsection*{About machines that write}

You may use a tool to look up a term, the way you would use a dictionary. You may not replace an interview, a test, or a file with a generated text and sign it. If you cannot point at the person, the screen, or the line you claim, the exercise is not done.


\chapter{Foreword}

Josef Müller-Brockmann wrote \emph{Grid Systems in Graphic Design} so that alignment would be a decision and not a habit. Massimo Vignelli's \emph{The Vignelli Canon} is the short version of the same stubbornness. You will draw lines on other people's screens. The books do not draw them.

\begin{center}
\includegraphics[width=0.78\textwidth,height=0.32\textheight,keepaspectratio]{../../resources/intros/UI-02-rules.jpg}\\[0.4em]
{\footnotesize Picture: AlanKitching, CC BY 4.0}
\end{center}

Recommended watching. The videos are context. They are not the deliverable.
\begin{itemize}
\item UI - Introduction \& Quest n°1, 01Edu \url{https://www.youtube.com/watch?v=UbE6WMeuvQg&list=PLHyAJ_GrRtf-_aBft2jQL4-QAunn2bLvu}
\end{itemize}

\chapter{Exercise 00: UI III - Ex 1 - Grids}
\exnumber{00}
\extitle{UI III - Ex 1 - Grids}
\exfiles{ex00/Figma file, named Name\_FirstName\_DeliverableName\_Date\_VersionNumber}
\exauthorize{Those named in the exercise}
\makeheaderfiles

\textbf{Context:}

Welcome to the 3rd quest of UI.

In this series of exercises, you'll be asked to open your eyes and collect graphic elements from various websites. You'll see grids across screens, multi state buttons, how to make a design accessible for all, breadcrumbs, toggle buttons, radio buttons, calendar, time picker with 2 main goals:

\begin{itemize}
\item Enrich your graphic general knowledge
\item Detect consistency, that helps users understand a digital product at the first glance.
\end{itemize}

\textbf{Instructions:}

Open these 5 images on Figma and draw lines and grids every time you find aligned elements.

\textbf{Deliverables:}

\begin{itemize}
\item A Figma file including the screens with lines on top to show grids. One page per screen or frame.
\item Upload the Figma file on Github, labeled as follows: "Name\_FirstName\_DeliverableName\_Date\_VersionNumber". For example, it can be labeled as "Doe\_John\_MidFiPrototype\_05242024\_V1".
\item Add a title within the file.
\item Pay attention to the global aspect of the file. It must be clear, simple and easy to read. You can get inspiration on canva.com but don't overload your design with too much details!
\end{itemize}

\textbf{Resources:}

\begin{itemize}
\item Video tutorial on Figma \url{https://www.youtube.com/watch?v=FTFaQWZBqQ8}
\item Beginning graphic design \url{https://edu.gcfglobal.org/en/beginning-graphic-design/}
\item Gareth graphic studio \url{https://www.youtube.com/c/GarethDavidStudio}
\item 1st screen example \url{https://www.figma.com/file/KOzd8yZuJvGLmZLuXjBgDg/UI-III-Ex-1.1} 
\end{itemize}


\chapter{Exercise 01: UI III - Ex 2 - Multi state}
\exnumber{01}
\extitle{UI III - Ex 2 - Multi state}
\exfiles{ex01/Figma file, named Name\_FirstName\_DeliverableName\_Date\_VersionNumber}
\exauthorize{Those named in the exercise}
\makeheaderfiles

\textbf{Instructions}:

Go to Sony.com \url{http://sony.com} and observe the buttons and call to actions. When you hover or click on them, they change!

Make screenshots of the different states of buttons, menu titles, filters etc. Then, gather and organize them.

Find at least five combinations of multi state elements.

\textbf{Deliverables:}

\begin{itemize}
\item A Figma file with 5 combinations of multi state elements.
\item Upload the Figma file on Github, labeled as follows: "Name\_FirstName\_DeliverableName\_Date\_VersionNumber". For example, it can be labeled as "Doe\_John\_MidFiPrototype\_05242024\_V1".
\item Add a title within the file.
\item Pay attention to the global aspect of the file. It must be clear, simple and easy to read. You can get inspiration on canva.com but don't overload your design with too much details!
\end{itemize}

\textbf{Resources:}

\begin{itemize}
\item Video tutorial on Figma \url{https://www.youtube.com/watch?v=FTFaQWZBqQ8}
\item Beginning graphic design \url{https://edu.gcfglobal.org/en/beginning-graphic-design/}
\item Gareth graphic studio \url{https://www.youtube.com/c/GarethDavidStudio}
\item Example here with Nintendo.com \url{https://www.figma.com/file/ahSnWRZeKBO62oJDiXltxY/UI-III---Ex-2}
\end{itemize}

/!\textbackslash{} The screenshots date back to nov 2021. They may differ from the current version of Nintendo.com \url{http://Nintendo.com} but the rules are the same


\chapter{Exercise 02: UI III - Ex 3 - Accessibility (website)}
\exnumber{02}
\extitle{UI III - Ex 3 - Accessibility (website)}
\exfiles{ex02/Figma file, named Name\_FirstName\_DeliverableName\_Date\_VersionNumber}
\exauthorize{Those named in the exercise}
\makeheaderfiles

\textbf{Instructions:}

Choose a website you find particularly accessible. Find at least 3 concrete examples of accessible measures and explain how they ease the reading or perception of visual content.

\textbf{Deliverables:}

\begin{itemize}
\item A 3-page Figma file with 3 different elements from a website (it can be from 3 different websites) that are accessible. The three measures must be different.
\item Upload the Figma file on Github, labeled as follows: "Name\_FirstName\_DeliverableName\_Date\_VersionNumber". For example, it can be labeled as "Doe\_John\_MidFiPrototype\_05242024\_V1".
\item Add a title within the file.
\item Pay attention to the global aspect of the file. It must be clear, simple and easy to read. You can get inspiration on canva.com but don't overload your design with too much details!
\end{itemize}

\textbf{Resources:}

\begin{itemize}
\item What accessibility is and why it's so important \url{https://uxdesign.cc/what-accessibility-is-and-why-its-so-important-9c56e033ff26}
\item Web Content Accessibility Guidelines (WCAG) \url{https://www.w3.org/WAI/standards-guidelines/wcag/}
\item Accessibility testing  ---  W3C Wik \url{https://www.w3.org/wiki/Accessibility_testing}i
\item UX Myth: Accessibility is expensive and difficult \url{https://uxmyths.com/post/654091803/myth-5-accessibility-is-expensive-and-difficult}
\item Salesforce UX - 7 tips every designer needs to know about accessibility \url{https://medium.com/salesforce-ux/7-things-every-designer-needs-to-know-about-accessibility-64f105f0881b}
\item Why prioritze web accessibility \url{https://blog.hubspot.com/service/why-prioritize-web-accessibility}
\end{itemize}

\textbf{Criteria:}

\begin{itemize}
\item Examples of measures
\begin{itemize}
\item Color contrast
\item Highlighted or enlarged text when hovering
\item Links color that is different from the body text's
\item Breadcrumbs to indicate where you are in terms of navigation
\item In forms, the title of the case should be above the text zone
\item Error states
\item etc.
\end{itemize}
\end{itemize}


\chapter{Exercise 03: UI III - Ex 4 - Accessibility (App)}
\exnumber{03}
\extitle{UI III - Ex 4 - Accessibility (App)}
\exfiles{ex03/Figma file, named Name\_FirstName\_DeliverableName\_Date\_VersionNumber}
\exauthorize{Those named in the exercise}
\makeheaderfiles

\textbf{Instructions:}

Choose an app you find particularly accessible. Find at least 3 concrete examples of accessible measures and explain how they ease the reading or perception of visual content.

\textbf{Deliverables:}

\begin{itemize}
\item A 3-page Figma file with 3 different elements from an app (it can be from 3 different apps) that are accessible. The three measures must be different.
\item Upload the Figma file on Github, labeled as follows: "Name\_FirstName\_DeliverableName\_Date\_VersionNumber". For example, it can be labeled as "Doe\_John\_MidFiPrototype\_05242024\_V1".
\item Add a title within the file.
\item Pay attention to the global aspect of the file. It must be clear, simple and easy to read. You can get inspiration on canva.com but don't overload your design with too much details!
\end{itemize}

\textbf{Resources:}

\begin{itemize}
\item What accessibility is and why it's so important \url{https://uxdesign.cc/what-accessibility-is-and-why-its-so-important-9c56e033ff26}
\item Web Content Accessibility Guidelines (WCAG) \url{https://www.w3.org/WAI/standards-guidelines/wcag/}
\item Accessibility testing  ---  W3C Wik \url{https://www.w3.org/wiki/Accessibility_testing}i
\item UX Myth: Accessibility is expensive and difficult \url{https://uxmyths.com/post/654091803/myth-5-accessibility-is-expensive-and-difficult}
\item Salesforce UX - 7 tips every designer needs to know about accessibility \url{https://medium.com/salesforce-ux/7-things-every-designer-needs-to-know-about-accessibility-64f105f0881b}
\end{itemize}

\textbf{Criteria:}

\begin{itemize}
\item Examples of measures
\begin{itemize}
\item Color contrast
\item Highlighted or enlarged text when hovering
\item Links color that is different from the body text's
\item Breadcrumbs to indicate where you are in terms of navigation
\item In forms, the title of the case should be above the text zone
\item Error states
\item etc.
\end{itemize}
\end{itemize}


\chapter{Exercise 04: UI III - Ex 5 - Breadbcrumbs}
\exnumber{04}
\extitle{UI III - Ex 5 - Breadbcrumbs}
\exfiles{ex04/Figma file, named Name\_FirstName\_DeliverableName\_Date\_VersionNumber}
\exauthorize{Those named in the exercise}
\makeheaderfiles

\textbf{Instructions:}

Browse the internet and find at least 5 different examples of breadcrumbs from websites or apps.

\textbf{Deliverable:}

\begin{itemize}
\item A Figma file gathering 5 screenshots of breadcrumbs.
\item Upload the Figma file on Github, labeled as follows: "Name\_FirstName\_DeliverableName\_Date\_VersionNumber". For example, it can be labeled as "Doe\_John\_MidFiPrototype\_05242024\_V1".
\item Add a title within the file.
\item Pay attention to the global aspect of the file. It must be clear, simple and easy to read. You can get inspiration on canva.com but don't overload your design with too much details!
\end{itemize}

\textbf{Resources:}

\begin{itemize}
\item What are breadcrumbs? \url{https://www.seoptimer.com/blog/breadcrumbs-website/}
\end{itemize}


\chapter{Exercise 05: UI III - Ex 6 - Toggle buttons}
\exnumber{05}
\extitle{UI III - Ex 6 - Toggle buttons}
\exfiles{ex05/Figma file, named Name\_FirstName\_DeliverableName\_Date\_VersionNumber}
\exauthorize{Those named in the exercise}
\makeheaderfiles

\textbf{Instructions:}

Browse the internet and find at least 5 different examples of toggle buttons from websites or apps.

\textbf{Deliverable:}

\begin{itemize}
\item A Figma file gathering 5 screenshots of toggle buttons.
\item Upload the Figma file on Github, labeled as follows: "Name\_FirstName\_DeliverableName\_Date\_VersionNumber". For example, it can be labeled as "Doe\_John\_MidFiPrototype\_05242024\_V1".
\item Add a title within the file.
\item Pay attention to the global aspect of the file. It must be clear, simple and easy to read. You can get inspiration on canva.com but don't overload your design with too much details!
\end{itemize}

\textbf{Resources:}

\begin{itemize}
\item Toggle Switch - Tips \url{https://uxplanet.org/toggle-switch-5-simple-design-tips-for-better-design-b4046eff4a2f}
\end{itemize}


\chapter{Exercise 06: UI III - Ex 7 - Radio Buttons}
\exnumber{06}
\extitle{UI III - Ex 7 - Radio Buttons}
\exfiles{ex06/Figma file, named Name\_FirstName\_DeliverableName\_Date\_VersionNumber}
\exauthorize{Those named in the exercise}
\makeheaderfiles

\textbf{Instructions:}

Browse the internet and find at least 5 different examples of radio buttons from websites or apps.

\textbf{Deliverable:}

\begin{itemize}
\item A Figma file gathering 5 screenshots of radio buttons.
\item Upload the Figma file on Github, labeled as follows: "Name\_FirstName\_DeliverableName\_Date\_VersionNumber". For example, it can be labeled as "Doe\_John\_MidFiPrototype\_05242024\_V1".
\item Add a title within the file.
\item Pay attention to the global aspect of the file. It must be clear, simple and easy to read. You can get inspiration on canva.com but don't overload your design with too much details!
\end{itemize}

\textbf{Resources:}

\begin{itemize}
\item What is a radio button? \url{https://www.justinmind.com/blog/radio-button-design-examples/}
\end{itemize}


\chapter{Exercise 07: UI III - Ex 8 - Calendars}
\exnumber{07}
\extitle{UI III - Ex 8 - Calendars}
\exfiles{ex07/Figma file, named Name\_FirstName\_DeliverableName\_Date\_VersionNumber}
\exauthorize{Those named in the exercise}
\makeheaderfiles

\textbf{Instructions:}

Browse the internet and find at least 5 different examples of calendars from websites or apps.

\textbf{Deliverable:}

\begin{itemize}
\item A Figma file gathering 5 screenshots of calendars.
\item Upload the Figma file on Github, labeled as follows: "Name\_FirstName\_DeliverableName\_Date\_VersionNumber". For example, it can be labeled as "Doe\_John\_MidFiPrototype\_05242024\_V1".
\item Add a title within the file.
\item Pay attention to the global aspect of the file. It must be clear, simple and easy to read. You can get inspiration on canva.com but don't overload your design with too much details!
\end{itemize}


\chapter{Exercise 08: UI III - Ex 9 - Time pickers}
\exnumber{08}
\extitle{UI III - Ex 9 - Time pickers}
\exfiles{ex08/Figma file, named Name\_FirstName\_DeliverableName\_Date\_VersionNumber}
\exauthorize{Those named in the exercise}
\makeheaderfiles

\textbf{Instructions:}

Browse the internet and find at least 5 different examples of time pickers from websites or apps.

\textbf{Deliverable:}

\begin{itemize}
\item A Figma file gathering 5 screenshots of time pickers.
\item Upload the Figma file on Github, labeled as follows: "Name\_FirstName\_DeliverableName\_Date\_VersionNumber". For example, it can be labeled as "Doe\_John\_MidFiPrototype\_05242024\_V1".
\item Add a title within the file.
\item Pay attention to the global aspect of the file. It must be clear, simple and easy to read. You can get inspiration on canva.com but don't overload your design with too much details!
\end{itemize}


\chapter{Exercise 09: UI III - Ex 10 - Micro-interactions}
\exnumber{09}
\extitle{UI III - Ex 10 - Micro-interactions}
\exfiles{ex09/Figma file, named Name\_FirstName\_DeliverableName\_Date\_VersionNumber}
\exauthorize{Those named in the exercise}
\makeheaderfiles

\textbf{Instructions:}

Browse the internet and find at least 5 different examples of feedback micro-interactions or micro-animations on websites or apps.

Be careful! Feedback is helpful to provide information to the user, not just showing a catchy design. You need to make the difference.

\textbf{Deliverable:}

\begin{itemize}
\item A Figma file gathering 5 examples of \textbf{feedback} micro-interactions or animations. In each example, show the different versions of the element.
\item Upload the Figma file on Github, labeled as follows: "Name\_FirstName\_DeliverableName\_Date\_VersionNumber". For example, it can be labeled as "Doe\_John\_MidFiPrototype\_05242024\_V1".
\item Add a title within the file.
\item Pay attention to the global aspect of the file. It must be clear, simple and easy to read. You can get inspiration on canva.com but don't overload your design with too much details!
\end{itemize}

\textbf{Resources:}

\begin{itemize}
\item Why use micro-interactions in your design \url{https://hike.one/update/why-use-micro-animations-in-your-design}
\item Micro-interactions: why, when,  and how to use them to boost the UX \url{https://uxdesign.cc/micro-interactions-why-when-and-how-to-use-them-to-boost-the-ux-17094b3baaa0}
\end{itemize}

\textbf{Tips:}

\begin{itemize}
\item Feedback is the situation when the app specifically sends a piece of information to the user, like: 
\item Payment accepted
\item Refreshed feed
\item Information saved
\item etc.
\end{itemize}


\chapter{Exercise 10: UI III - Ex 11 - Consistency}
\exnumber{10}
\extitle{UI III - Ex 11 - Consistency}
\exfiles{ex10/Figma file, named Name\_FirstName\_DeliverableName\_Date\_VersionNumber}
\exauthorize{Those named in the exercise}
\makeheaderfiles

\textbf{Instructions:}

Pick a website or an app of your choice*. Highlight 5 signs of consistency across the pages.

*Airbnb, Google, Apple, Amazon, Facebook, Instagram and WhatsApp are not accepted.

\textbf{Deliverable:}

\begin{itemize}
\item A Figma file gathering at least 5 signs of consistency.
\item Upload the Figma file on Github, labeled as follows: "Name\_FirstName\_DeliverableName\_Date\_VersionNumber". For example, it can be labeled as "Doe\_John\_MidFiPrototype\_05242024\_V1".
\item Add a title within the file.
\item Pay attention to the global aspect of the file. It must be clear, simple and easy to read. You can get inspiration on canva.com but don't overload your design with too much details!
\end{itemize}

\textbf{Resources:}

\begin{itemize}
\item 7 tips for designing consistency \url{https://designshack.net/articles/graphics/7-tips-for-designing-consistency/}
\item The value of consistent design \url{https://www.invisionapp.com/inside-design/consistent-design/}
\item Consistency and Standards - Examples of Nielsens's Design Heuristics \url{https://medium.com/@gregoralbrecht/consistency-and-standards-nielsens-design-heuristic-explained-ac91f450fd8}
\item The importance of consistency in design work \url{https://yesimadesigner.com/the-importance-of-consistancy-in-design-work/}
\item Examples of signs of consistency: 
\begin{itemize}
\item colors
\item same shaped buttons
\item alignment
\item consistent product pages (with identical frameworks)
\item etc.
\end{itemize}
\end{itemize}


\end{document}
